% Original figure: the 20 standard amino acids grouped by side-chain property, and the peptide
% bond (plan-report 3.1). Standard biochemistry; grouping is one common convention (histidine is
% listed with the basic residues although it is only partly charged at neutral pH).
\begin{minipage}[t]{8.3cm}
\begin{tikzpicture}[fig, baseline=(current bounding box.north)]
  % ---------------- (a) groups
  \node[title, anchor=west] at (0,2.65) {(a) Twenty standard amino acids by side chain};
  \newcommand{\aagroup}[5]{% y, colour, label, letters, names
    \node[font=\scriptsize\bfseries, anchor=west, text width=2.45cm, align=left] at (0,#1) {#3};
    \foreach \a [count=\k] in {#4} {
      \node[residue, minimum size=4.6mm, fill=#2] at (2.35+0.52*\k,#1) {\a};
    }
    \node[font=\scriptsize, anchor=west, text=oiGrey] at (5.6,#1) {#5};
  }
  \aagroup{2.15}{oiYellow!35}{hydrophobic}{A,V,L,I,M}{Ala Val Leu Ile Met}
  \aagroup{1.6}{oiOrange!30}{aromatic}{F,W,Y}{Phe Trp Tyr}
  \aagroup{1.05}{oiSky!30}{polar}{S,T,N,Q}{Ser Thr Asn Gln}
  \aagroup{0.5}{oiBlue!25}{basic ($+$)}{K,R,H}{Lys Arg His}
  \aagroup{-0.05}{oiVermillion!25}{acidic ($-$)}{D,E}{Asp Glu}
  \aagroup{-0.6}{oiGreen!25}{special}{G,P,C}{Gly Pro Cys}
  \node[note, text width=7.6cm, align=left, anchor=north west] at (0,-0.95)
       {One-letter codes are used throughout. A substitution replaces one residue by any of the
        other 19, so each position has 19 possible single substitutions.};

\end{tikzpicture}
\end{minipage}\hfill
\begin{minipage}[t]{6.9cm}
\begin{tikzpicture}[fig, baseline=(current bounding box.north)]
  \node[title, anchor=west] at (0,2.65) {(b) Peptide bond};
  \node[anchor=north west] (pb) at (0,2.25) {\scriptsize
    \chemfig{H_2N-C(-[2]R_1)(-[6]H)-C(=[2]O)-[0,,,,draw=oiVermillion,line width=1.6pt]N(-[6]H)-C(-[2]R_2)(-[6]H)-COOH}};
  \node[font=\scriptsize, text=oiVermillion, anchor=north] at ($(pb.south)+(0.1,0.05)$) {peptide bond (C--N)};
  \node[note, text width=6.5cm, align=left, anchor=north west] at (0,-0.05)
       {A peptide is a chain of residues joined N to C; a 28-residue tile is such a chain. The side
        chains $R$ differ between amino acids and carry most of the chemical difference a
        substitution introduces.};
\end{tikzpicture}
\end{minipage}
